\label{chap:83-12}

\section{The Hensel--Witt Transduction Kernel}
\label{p12-83-c12-s1:chap:nthw}
\index{Hensel--Witt Transduction Kernel}

\readerquestion{Where do the cylinders, the monodromy, the
five regions, the hexads, the octads and the physical readers actually intersect?}

\subsection{Hensel–Witt Transduction Kernel: microregions of \(\pi\), holonomy \(\Gamma_9\), dodecaphasic register and hexad–octad transport}

The NTHW is the data and maps package
\[
\NTHW=
\bigl(
\mathscr R_\pi^{\rm micro},
\Gnine,
\Cdoce,
K,
\mathcal W_{12\to24},
\mathcal I_{6\to8}
\bigr),
\]
with the following compatibilities:
\begin{enumerate}
  \item the Hensel regions share a quotient word and retain
  distinct memory;
  \item \(\Gnine\) transports orientation and produces monodromy;
  \item \(\Cdoce\) organizes the signed observable and the seal \(K\);
  \item the Witt reading converts supports of the same boundary into hexads and
  stars;
  \item the extension \(W_{12}\to W_{24}\) produces the
  hexad--octad interface;
  \item the Archimedean, exceptional and dimensional readers receive
  invariants of the same state.
\end{enumerate}

\begin{figure}[p]
\centering
\resizebox{\textwidth}{!}{\input{figures/fig_nucleo_hensel_witt}}
\caption{The NTHW as a typed interface. The central mathematical task is to make
the diagram commute, not to accumulate coincidences around a name.}
\label{p12-83-c12-s1:fig:nthw}
\end{figure}

NTHW gives a formal name to the place previously described as the “Rosetta
Stone.” The substitution is conceptual: a stone translates through resemblance; a
transduction specifies input, output, preserved datum and loss of
information.

\subsection{From the ternary code to the Witt system}

The reversible APP boundary selects, with declared gauges, a matrix
\(A_W\) over \(\FF_3\) satisfying
\[
A_WA_W^T=2I\pmod3.
\]
The graphic code
\[
C_W=\{(q,qA_W):q\in\FF_3^6\}
\]
is the extended ternary Golay code. Its weight enumerator is
\[
1+264y^6+440y^9+24y^{12}.
\]
The \(264\) words of weight six, identified up to sign, yield \(132\)
supports. Every five-point subset belongs to a unique hexad:
\[
W_{12}=S(5,6,12).
\]

The field of \(495\) star involutions generates a group of order
\(95\,040\), recognized as \(M_{12}\). An isolated star does not generate the
group. The ordered HMT frame has trivial stabilizer and therefore fixes a
gauge within the action.

\subsection{Hexad, octad and the corridor \texorpdfstring{\(90/120\)}{90/120}}

Let \(H\) be a hexad and \(O_H\) an octad containing it. Marking one point
and a pair of points of the hexad yields
\[
F_6(H)=H\times\binom H2,
\qquad
|F_6|=6\binom62=90.
\]
If the marked point can traverse the octad:
\[
F_8(H)=O_H\times\binom H2,
\qquad
|F_8|=8\binom62=120.
\]
The inclusion \(H\hookrightarrow O_H\) induces \(F_6\hookrightarrow F_8\).
Its normalized defect is
\[
\boxed{
\frac{90-120}{24\binom62}
=\frac{6-8}{24}
=-\frac1{12}.}
\]
This combinatorial theorem precedes Barbero’s angular reading. The
two objects share the interface \(6\to8\), but are not identical.

\subsection{\(4\) rutas hacia -1/12}

The construction contains four realizations that must be presented together, not
counted as four independent proofs of a single provenance.

\subsubsection{Elementary Operator of the Dodecaphasic Cycle and Update of the Signed Phase}

One tooth of \(\Cdoce\) measures \(30^\circ\). Then
\[
\frac{30}{120}-\frac{30}{90}
=\frac14-\frac13
=-\frac1{12}.
\]
It is the angular shadow of the pattern \(90/120\).

\subsubsection{Gram matrix of the surviving sections and positivity criterion}

In the certified jump \(60\to81\), two fibers have sizes \(0\) and
\(12\). The Gram operator is
\[
K_{60,81}=\operatorname{diag}(0,12).
\]
On the live subspace, its pseudoinverse is \(1/12\); with the negative
boundary orientation,
\[
E_{\rm surv}=-K_{60,81}^{+}
\]
yields \(-1/12\).

\subsubsection{Triadic scale extractor}

The realization via a triadic scale extractor is the second-difference construction
developed in full in
\cref{sec:c12-segunda-diferencia-triadica}. This first occurrence establishes its
role within the inventory of realizations; the indicated residence preserves
the definitions of \(T_L\), \(a(L)\), \(A_1(L)\), \(C(L)\), the independence
of the residue from \(L\), and the exact status of the normalization.

\subsubsection{Modular corridor between dodecaphasic cards and conservation of the phase invariant}

The eta function satisfies
\[
\frac{\eta(\tau)}{\eta(3\tau)}
=q^{-1/12}\prod_{n\ge1}(1+q^n+q^{2n})^{-1}.
\]
The exponent comes from \(1/24-3/24=-1/12\). Taking twelve copies:
\[
t_3(q)=\left(\frac{\eta(\tau)}{\eta(3\tau)}\right)^{12}
=q^{-1}-12+54q-76q^2-243q^3+\cdots.
\]
The coefficient \(54\) coincides with the APP defect. The corridor
\[
j=\frac{(t_3+27)(t_3+243)^3}{t_3^3}
\]
connects afterwards with the classical modular branch.

\begin{exactbox}[title={What has been proved in the corridor}]
The APP defect \(54\), the eta exponent \(-1/12\), the
hexad--octad defect, and the modular identities are exact once \(t_3\) has been fixed. The
strong assertion is that all of them are natural images of a single
TPK operator. That naturality must be proved through the NTHW diagram; it is not
deduced merely from repeating the same integers.
\end{exactbox}

\subsection{Prime-power defect and number operator}

For a prime power \(p^m\), the complete APP defect generalizes to
\[
\Delta_{p^m}^{\rm full}
=\frac{m(p-1)}2p^{2m-1}.
\]
In particular,
\[
\Delta_3=3,
\qquad
\Delta_9=54,
\qquad
\Delta_{27}=729.
\]
Upon normalization,
\[
a_{p,m}
=\frac{2\Delta_{p^m}^{\rm full}}{(p-1)p^{2m-1}}
=m.
\]
With \(x=p^{-s}\):
\[
\sum_{m\ge1}a_{p,m}x^m
=\sum_{m\ge1}mx^m
=\frac{x}{(1-x)^2}
=x\frac{\dd}{\dd x}\frac1{1-x}.
\]
The factor \((1-x)^{-1}\) is the local Euler factor; differentiation inserts the
number operator. Thus, an exact bridge appears:
\[
\boxed{
\text{tower of APP defects in }p^m
\longleftrightarrow
\text{insertion of the number operator into the local Euler factor}.}
\]

\subsection{Joint naturality conditions for the Archimedean, exceptional, and dimensional transduction}

The NTHW will be a complete transduction theorem when there is a family of
invariants \(\mathcal I\) and readers
\[
\mathfrak R,\quad\mathfrak E,\quad\mathfrak D
\]
such that
\[
\mathcal I
=\mathcal I_{\rm arq}\circ\mathfrak R
=\mathcal I_{\rm exc}\circ\mathfrak E
=\mathcal I_{\rm dim}\circ\mathfrak D,
\]
and when the maps of region, seal, incidence, and action are derived from the same
transport. The finite proofs already provide candidates: monodromy,
support \(B_K\), borrowing pattern, two-way parity, and counts \(90/120\).

\begin{programbox}[title={Executable test 108--144--90--120}]
Label each of the \(1008\) routes of \(\pi\) by CT108 orbit and by
Witt completion; compute the fibers of the map; test whether the four
regions of multiplicity \(144\) realize factor \(8/6\) on blocks
of \(108\); and verify whether the region of \(432\) transports the monodromy. A
single counterexample to label compatibility would refute this bridge in
its present form.
\end{programbox}

\begin{takeawaybox}
The NTHW is the true center of the account because there the same boundary may be read
as cylinder, monodromy, seal, incidence, and action. Its power does not consist
in many numbers resembling one another, but in several finite maps already
sharing supports and exact defects. The next task is to prove the
joint naturality of those charts.
\end{takeawaybox}


\section{Marked inclusion of the hexad-octad incidence and preservation of its multiplicities}

With the dodecad and the alignment of the
\cref{p12-83-c17-s1:sec:w12-w24-elevacion} fixed, let \(H\) be a hexad in the HMT coordinates,
let \(\ell(H)\subset D\) be its image, and let
\(O_H=\iota_D(\ell(H))\) be its unique lift. Marking a point
and a pair of points of the residual hexad yields
\[
F_6(H)=\ell(H)\times\binom{\ell(H)}2,
\qquad
|F_6|=6\binom62=90.
\]
If the marked point can traverse the octad:
\[
F_8(H)=O_H\times\binom{\ell(H)}2,
\qquad
|F_8|=8\binom62=120.
\]
The inclusion \(\ell(H)\hookrightarrow O_H\) induces the incidence morphism
\[
 \mathcal I_{6\to8}:F_6(H)\hookrightarrow F_8(H).
\]
Normalizing the difference by the chosen
factor \(24\binom62\) yields
\[
\boxed{
\frac{90-120}{24\binom62}
=\frac{6-8}{24}
=-\frac1{12}.}
\]
The equality is a normalized defect of the combinatorial interface \(6\to8\). The factor \(24\binom62\) forms part of the definition of this normalization; the inclusion alone determines the difference \(-30\). The operator \(\mathcal I_{6\to8}\) has marked incidence configurations as its domain; it is neither the transverse extractor of the dodecaphase state nor the family of angular moments.

\section{Classification by domain of the \(5\) algebraic realizations of \(-1/12\)}

In addition to the normalized defect of the hexad--octad inclusion,
five operator realizations of the same rational number appear. Their domains are
specified separately; equality of values does not identify the
operators without an additional intertwiner.

\subsection{Difference of dodecaphasic periodicities}

One tooth of \(\Cdoce\) measures \(30^\circ\). Then
\[
\frac{30}{120}-\frac{30}{90}
=\frac14-\frac13
=-\frac1{12}.
\]
It is the oriented difference between the periods of \(90^\circ\) and
\(120^\circ\) in this realization.

\subsection{Pseudoinverse of the Gram Operator on the Live Component}

In the exact jump \(60\to81\), two fibers have sizes \(0\) and
\(12\). The Gram operator is
\[
K_{60,81}=\operatorname{diag}(0,12).
\]
On the nonzero component of its range, the pseudoinverse is \(1/12\); with the negative
boundary orientation,
\[
E_{\rm surv}=-K_{60,81}^{+}
\]
yields \(-1/12\).

\subsection{2nd difference between triadic scales}
\label{sec:c12-segunda-diferencia-triadica}

Define
\[
T_L=\sum_{n=1}^{L-1}n(L-n)=\frac{L(L^2-1)}6,
\qquad
a(L)=-\frac{T_L}{L^2}=-\frac L6+\frac1{6L}.
\]
The first scale subtraction is
\[
A_1(L)=a(L)-\frac{a(3L)}3=\frac4{27L},
\]
and the second
\[
C(L)=LA_1(L)-\frac{3L}{3}A_1(3L)=\frac8{81}.
\]
The residue \(8/81\) is independent of \(L\). To obtain
\(-1/12\), the declared normalization is applied
\[
R(L)=-\frac{27}{32}C(L)=-\frac1{12}.
\]
The calculation derives \(8/81\); it does not by itself derive the factor \(-27/32\).

\subsection{Composition of modular transformations and preservation of the signed state}

The eta function satisfies
\[
\frac{\eta(\tau)}{\eta(3\tau)}
=q^{-1/12}\prod_{n\ge1}(1+q^n+q^{2n})^{-1}.
\]
The exponent comes from \(1/24-3/24=-1/12\). Taking twelve copies:
\[
t_3(q)=\left(\frac{\eta(\tau)}{\eta(3\tau)}\right)^{12}
=q^{-1}-12+54q-76q^2-243q^3+\cdots.
\]
The coefficient \(54\) coincides with the APP defect\@. The rational identity
\[
j=\frac{(t_3+27)(t_3+243)^3}{t_3^3}
\]
establishes the connection with the classical modular branch.

\subsection{Defect of cyclic difference projectors}

Let \(S\) be the cyclic shift on \(\QQ^{12}\) and define
\[
 D_r=I-S^r.
\]
The kernel of \(D_r\) consists of the vectors constant on each orbit
of \(S^r\); therefore,
\[
 \dim\ker D_r=\gcd(12,r).
\]
Let \(\Pi_{\ker D_3}\) and \(\Pi_{\ker D_4}\) be the rational projectors onto
\(\ker D_3\) and \(\ker D_4\), and let
\(\tau(T)=\operatorname{Tr}(T)/12\) be the normalized trace. Then
\[
 \boxed{
 \tau(\Pi_{\ker D_3}-\Pi_{\ker D_4})
 =\frac{\operatorname{rank}\Pi_{\ker D_3}
       -\operatorname{rank}\Pi_{\ker D_4}}{12}
 =\frac{3-4}{12}
 =-\frac1{12}.}
\]
This realization expresses the rational as the transverse defect of steps
three and four in the twelve-phase cycle.

\begin{remark}[Domains of the realizations]
The APP defect \(54\), the eta exponent \(-1/12\), the combinatorial
defect of the hexad--octad lift, and the
modular identities once \(t_3\) has been fixed are exact. The value
of the pseudoinverse, the normalized extractor, and the projector defect are likewise exact.
Each equality preserves its domain and its corresponding map.
\end{remark}
